\documentclass[a4paper,12pt]{article}
\usepackage[utf8]{inputenc} 
\usepackage{amsmath}
\usepackage[a4paper, margin=0.3in,top=0.1in, bottom=0.3in]{geometry} % Adjust the margin here
\usepackage{multirow}
\usepackage{circuitikz}
\usepackage{nopageno}
\usepackage{graphicx} % Add this line to include graphics
\setlength{\parindent}{0pt} % Set global no indent
    
\title{{\large Department of Physics, FNSPE, CTU in Prague} \\
                  02YELMA - Electricity and Magnetism \\ 
          {\Large \textbf{Final Exam - Solutions 2}}\vspace{-1cm}}
\date{}
\author{}

\begin{document}

\maketitle
\vspace{-0.9cm}

\noindent

\textbf{Q1:} (a) Energy is conserved in the circuit, so the total voltage drop across the circuit elements must equal the applied voltage $V=V_R+V_L+V_C=0$. The voltage drops are

\[
V_R = R I, \quad V_L = L\frac{dI}{dt}, \quad V_C = \frac{q}{C}
\]

The charge on the capacitor is related to the current by $I=dq/dt$. Thus,

\[
R I
+
L\dot{I}
+
\frac{q}{C} = 0 ~~\xrightarrow{d/dt}~~ 
R\dot{I}
+
L\ddot{I}
+
\frac{1}{C}I=0 \implies \boxed{
\ddot{I}
+
\frac{{R}}{L}\dot{I}
+
\frac{1}{LC}I
=
0
}
\]

(b) The governing equation describes a damped oscillation where the coefficient of the first derivative term determines the amount of damping: $2\beta = R/L$. To obtain the characteristic equation, we assume a solution of the form $I(t)=e^{\lambda t}$, which gives
\[
\lambda^2 + \frac{R}{L}\lambda + \frac{1}{LC} = 0 \implies \lambda = -\frac{R}{2L} \pm \sqrt{\left(\frac{R}{2L}\right)^2 - \frac{1}{LC}}.
\]
For underdamped oscillations, the discriminant must be negative, which gives the condition on $R$, $L$, and $C$:
\[
\boxed{
R<2\sqrt{L/C}
}
\]
For current to fall to half its initial value:
\[
\frac{I_\text{max}e^{-\beta t}}{I_\text{max}} = \frac{1}{2} \implies e^{-\beta t} = \frac{1}{2} \implies \boxed{t = \frac{\ln 2}{\beta} = \frac{2L\ln 2}{R}}
\]

(c) The total energy in the circuit is $U = U_L + U_C = \frac{1}{2}LI^2 + \frac{1}{2}CV^2$. If the capacitor is initially charged to $V_0$, then the initial energy is $U_0 = \frac{1}{2}CV_0^2$. When the current is at its maximum, the energy is stored entirely in the inductor, so $U = \frac{1}{2}LI_\text{max}^2$. At resonance, we know that only the resistor dissipates energy and current behaves as $I \sim I_0\sin(\omega_0 t)$. Therefore, the average energy dissopated in the resistor over one cycle is
\[
P_\text{avg} = \frac{1}{2}RI_0^2,~~T = \frac{2\pi}{\omega_0}  \implies \Delta E = P_\text{avg}T = \pi RI_0^2/\omega_0.
\]
Therefore, the quality factor is
\[
Q = 2\pi\frac{U}{\Delta E} = 2\pi\frac{\frac{1}{2}LI_0^2}{\pi RI_0^2/\omega_0} = \frac{\omega_0 L}{R} = \frac{1}{R}\sqrt{\frac{L}{C}} \implies \boxed{
Q = \frac{1}{R}\sqrt{\frac{L}{C}}
}
\]


\textbf{Q2:} (a) Because of cylindrical symmetry, the magnetic field can only have a angular component, $\vec B(\phi,t)=B(\phi,t)\,\hat \phi$. Consider an Amperian loop of radius \(a<s<b\) coaxial with the shells:
\[
B(\phi,t)\,(2\pi s)
=
\mu_0 I(t) \implies \boxed{\vec B(\phi,t)
=
\frac{\mu_0 I_0 \cos(\omega t)}{2\pi s}\,\hat \phi,
~~\text{for}~~a<s<b.
}
\]
In region \(s>b\), the net current enclosed by the Amperian loop is zero, so \(\vec B(\phi,t)=0\).

(b) The magnetic energy density is given by

\[
u_B=\frac{B^2}{2\mu_0}.
\]

The instantaneous magnetic energy stored per unit length is obtained by integrating over the vacuum region between the conductors:

\[
\frac{U(t)}{L}
=
\int_a^b
\frac{B^2(r,t)}{2\mu_0}\,(2\pi r\,dr)=\int_a^b
\frac{1}{2\mu_0}
\left(
\frac{\mu_0 I(t)}{2\pi r}
\right)^2
(2\pi r\,dr)=\frac{\mu_0 I^2(t)}{4\pi}
\ln\!\left(\frac{b}{a}\right)
.
\]

\[
\implies \boxed{
\frac{U(t)}{L}
=
\frac{\mu_0 I_0^2}{4\pi}
\cos^2(\omega t)
\ln\!\left(\frac{b}{a}\right)
}
\]

The inductance per unit length is defined through

\[
\frac{U}{L}
=
\frac{1}{2}\mathcal{L}I^2(t) \implies 
\boxed{
\mathcal{L}
=
\frac{\mu_0}{2\pi}
\ln\!\left(\frac{b}{a}\right)
}.
\]

\textbf{Q3:}The solenoid is at rest in frame $S'$ and produces the fields $\vec{E}'=0,~\vec{B}'=B_0\hat{z}'$. Frame $S'$ moves relative to the laboratory frame $S$ with velocity $\vec{v}=v\hat{x}$. Since the motion is along the $x$-direction, the field components parallel and perpendicular to $\vec v$ are $E'_{\parallel}=0,~E'_{\perp}=0$ and because $\vec B'$ points along $\hat z$, $B'_{\parallel}=0,~B'_{\perp}=B_0\hat z$. Using
\[
E_{\parallel}=E'_{\parallel},
~~
E_{\perp}=\gamma\left(E'_{\perp}-\vec v\times\vec B'\right) \qquad \implies \qquad E_{\parallel}=0,~~ \vec E
=
-\gamma(\vec v\times\vec B').
\]

Compute the cross product: $
\vec v\times\vec B'
=
(v\hat x)\times(B_0\hat z)
=
vB_0(\hat x\times\hat z)
=
-vB_0\hat y \implies \boxed{
\vec E
=
\gamma vB_0\,\hat y
}.
$

Using
\[
B_{\parallel}=B'_{\parallel},
~~
B_{\perp}
=
\gamma\left(
B'_{\perp}
+
\frac{1}{c^2}\vec v\times\vec E'
\right),~~\vec E'=0 \implies B_{\parallel}=0 \text{ and } \vec B
=
\gamma B_0\hat z.
\]


\[
\implies
\boxed{
\vec B
=
\gamma B_0\,\hat z
}.
\]


\textbf{Q4:} The conducting rod of length \(L\) moves along horizontal rails and closes a circuit containing two resistors \(R\) (rod resistance) and \(R_0\). The total resistance is $R_{\mathrm{tot}}=R+R_0$. Let \(x(t)\) denote the position of the rod, so the enclosed area is $A(t)=Lx(t)$. The magnetic flux through the loop is $\Phi_B(t)=B(t)Lx(t)$. Using Faraday's law, $\mathcal E=-\frac{d\Phi_B}{dt}$, we obtain

\[
\mathcal E
=
-\frac{d}{dt}\left(BLx\right)
=
-L\left(x\frac{dB}{dt}+B\frac{dx}{dt}\right) \implies \mathcal E=L B(t)\left(\alpha x-\dot x\right) \implies \boxed{
I(t)=\frac{L B_0 e^{-\alpha t}}{R+R_0}
\left(\alpha x-\dot x\right)
}.
\]
Direction follows Lenz's law:

- If \(\alpha x>\dot x\), then \(d\Phi_B/dt<0\): the magnetic flux through the loop decreases, so the induced current must create magnetic field in the \(+\hat z\) direction.
  Hence the current is \textbf{counterclockwise} (viewed from above).

- If \(\alpha x<\dot x\), the flux increases and the current reverses:
  \textbf{clockwise}.

Initially (\(v(0)=0\)):

\[
I(0)=
\frac{L B_0\alpha x_0}{R+R_0}>0,
\]

thus the initial current is counterclockwise.


(b) The rod experiences magnetic force $\vec F = I\,\vec L\times \vec B$. The rod carries current along \(\hat y\), therefore $\vec F = I L B\,\hat x$.

Substituting \(I\),

\[
F_x=
\frac{L^2 B^2}{R+R_0}
(\alpha x-\dot x) \implies m\ddot x
=
\frac{L^2 B_0^2 e^{-2\alpha t}}{R+R_0}
(\alpha x-\dot x) \implies \boxed{
m\ddot x
+
\frac{L^2 B_0^2 e^{-2\alpha t}}{R+R_0}\dot x
-
\frac{\alpha L^2 B_0^2 e^{-2\alpha t}}{R+R_0}x
=0
}
\]

(c) 
\[
I(0)=
\frac{L B_0\alpha x_0}{R+R_0} \implies P_R(0)=I(0)^2R
=
\frac{L^2B_0^2\alpha^2x_0^2\,R}
{(R+R_0)^2}, \qquad P_{R_0}(0)=I(0)^2R_0
=
\frac{L^2B_0^2\alpha^2x_0^2\,R_0}
{(R+R_0)^2}
\]

\[
\implies
\boxed{
P_{\mathrm{tot}}(0)
=
P_R(0)+P_{R_0}(0)
=
\frac{L^2B_0^2\alpha^2x_0^2}{R+R_0}
}
\]

\end{document}

